% ECONOMICS_PROBLEM: {"id": "D63-448", "title": "Fair chores with unequal time requirements", "jel_code": "D63", "source_id": "OP-0277"}
% FIRST_POSTED: 2026-10-10
% ATTEMPT_STATUS: partial
% ENTRY_KIND: research_attempt
\documentclass[11pt]{article}
\usepackage[a4paper,margin=28mm]{geometry}
\usepackage{amsmath,amssymb,amsthm}
\usepackage[T1]{fontenc}
\usepackage{lmodern,microtype}
\usepackage[colorlinks=true,allcolors=blue]{hyperref}
\newtheorem{theorem}{Theorem}
\newtheorem{corollary}[theorem]{Corollary}
\newcommand{\OPT}{\operatorname{OPT}}
\title{A sharp capacity boundary for budget-feasible EF1 chore allocations}
\author{Ji Ho Bae}
\date{10 October 2026}
\begin{document}
\maketitle
\begin{abstract}
Consider complete chore allocations with common processing times in a fixed
interval $[a,b]$ and a common time budget $B$, where $0<a\le b\le B$.
For any fixed number $n\ge2$ of agents, the best uniform EF1 factor is $1$
if $\lfloor B/a\rfloor=\lfloor B/b\rfloor$, and is unbounded otherwise.
The positive case follows from least-cost round robin. In the negative case,
a packing argument forces every feasible allocation to have the same unfair
shape, even with strictly positive identical disutilities. Thus arbitrarily
small differences in processing times can destroy every finite guarantee.
\end{abstract}

\section{Model and statement}

There are $n$ agents and a finite set $C$ of indivisible chores. Chore $c$
has a common processing time $t_c\in[a,b]$, and every agent has budget $B$.
Agent $i$ has additive nonnegative disutility
$d_i(S)=\sum_{c\in S}d_i(c)$. An allocation $A=(A_1,\ldots,A_n)$ is a
partition of $C$, and is feasible if $\sum_{c\in A_i}t_c\le B$ for all $i$.
We consider only instances admitting a feasible allocation.

For $\alpha\ge1$, call $A$ $\alpha$-EF1 if for every $i\ne j$ either
$A_i=\varnothing$, or some $c\in A_i$ satisfies
\[
                     d_i(A_i\setminus\{c\})\le\alpha d_i(A_j).
\]
The deleted chore belongs to the potentially envious agent. In this common-time,
common-budget setting every feasible bundle fits every agent, so this is also
the budget-aware full-bundle comparison proposed in catalogue entry D63-448.
Let $\alpha(A)$ be the infimum of the admissible factors, with
$\inf\varnothing=+\infty$, and define the uniform guarantee
\[
 G_n(a,b;B)=\sup_I\ \min_{A\text{ feasible}}\alpha(A),
\]
where $I$ ranges over the instances just specified. The bounds $a,b$ need not
both occur as processing times in an instance.

\begin{theorem}\label{thm:boundary}
For every $n\ge2$ and $0<a\le b\le B$,
\[
 G_n(a,b;B)=
 \begin{cases}
  1,&\lfloor B/a\rfloor=\lfloor B/b\rfloor,\\
  +\infty,&\lfloor B/a\rfloor>\lfloor B/b\rfloor.
 \end{cases}
\]
In the first case, least-cost round robin constructs a feasible EF1 allocation.
In the second, the guarantee remains unbounded for identical strictly positive
disutilities and only two processing times.
\end{theorem}

Elkind's problem in the Dagstuhl report~\cite[Section~4.1]{elkind} concerns
fairness with agent-dependent chore times and asks for suitable envy notions.
The displayed approximation criterion is the catalogue's formulation, not a
definition adopted by that report. A prior catalogue attempt~\cite{prior}
already gives an unbounded-factor example and proves the equal-duration
round-robin guarantee. Theorem~\ref{thm:boundary} determines the exact boundary
for every fixed interval of common times. It does not address the broader
heterogeneous-time axiomatization agenda. The budget model of
Elkind, Igarashi, and Teh~\cite[Section~2.1]{eit} permits an auxiliary
housekeeper; here every chore must be allocated to the $n$ agents.

\section{The capacity test is sufficient}

Write $k=\lfloor B/b\rfloor$. If $\lfloor B/a\rfloor=k$, every feasible
bundle has at most $k$ chores. Hence a feasible instance with $m$ chores
satisfies $m\le nk$.

Run round robin in a fixed agent order, each agent choosing a remaining chore
of minimum disutility on her turn. Bundle sizes differ by at most one, so each
agent receives at most $\lceil m/n\rceil\le k$ chores. Every resulting workload
is therefore at most $kb\le B$.

For completeness, fix agents $i\ne j$ with $A_i\ne\varnothing$ and remove
the last chore chosen by $i$. If $i$ precedes $j$ in the round, compare each
retained $r$th choice of $i$ with the $r$th choice of $j$. If $i$ follows $j$,
compare it with the $(r+1)$st choice of $j$. In either case the comparison
choice exists, was still available when $i$ chose, and differs for different
retained chores. Minimum-disutility selection and nonnegativity give
\[
                    d_i(A_i\setminus\{c\})\le d_i(A_j).
\]
Thus the allocation is EF1. Since factors are constrained to be at least one,
the guarantee is exactly one. With rational inputs, the construction takes a
polynomial number of exact comparisons.

\section{A capacity gap is sufficient for unboundedness}

Suppose $\lfloor B/a\rfloor>k=\lfloor B/b\rfloor$. Then $k\ge1$ and
\[
                       a\le s:=\frac{B}{k+1}<b.
\]
Fix $0<\varepsilon\le1$. Use $(n-1)k$ cheap chores, each with time $b$
and disutility $\varepsilon$, and $k+1$ expensive chores, each with time $s$
and disutility $1$. All agents have these same disutilities. The instance is
feasible: one agent receives all expensive chores, and every other agent
receives $k$ cheap chores.

There are $nk+1$ chores, so in every complete allocation some agent receives
at least $k+1$ chores. Every chore takes at least $s$, and $(k+1)s=B$.
A feasible bundle of that size consequently has exactly $k+1$ chores and
contains no cheap chore, since $b>s$. It therefore contains all expensive
chores. Each other bundle consists only of cheap chores and contains at most
$k$ of them. Completeness forces each of these $n-1$ bundles to contain
exactly $k$ cheap chores.

After any deletion, the expensive-bundle agent has disutility $k$ and compares
it with $k\varepsilon$. Thus every feasible allocation requires
$\alpha\ge1/\varepsilon$. This factor also suffices: the expensive-bundle
inequalities hold with equality, and each cheap-bundle agent has residual
$(k-1)\varepsilon$, no greater than any comparison disutility after
multiplication by $1/\varepsilon$. Hence every feasible allocation satisfies
\[
                           \alpha(A)=1/\varepsilon.
\]
Letting $\varepsilon$ decrease to zero proves the second case of
Theorem~\ref{thm:boundary}. All data in the construction are rational whenever
$a,b,B,\varepsilon$ are rational.

\begin{corollary}
For every fixed $n\ge2$ and every $R>1$, no finite factor is guaranteed over
feasible common-time instances with a nonempty chore set, equal budgets, and
$\max_c t_c/\min_c t_c\le R$, even with identical strictly positive
disutilities. At $R=1$, factor one is always attainable.
\end{corollary}
\begin{proof}
Take $B=2$, $a=1$, and $b=1+\eta$ with
$0<\eta<\min\{R-1,1\}$. The two capacities are two and one, so the negative
case applies. When $R=1$, all processing times are equal and the capacities
coincide.
\end{proof}


\section*{Status, verification, and attribution}
\noindent\textbf{Author: Ji Ho Bae.}

The fixed-interval theorem above is proved in full. This is a partial result
for D63-448 as a whole: heterogeneous processing times, unequal budgets,
and the broader normative axiomatization are not classified here.
No numeric completion estimate is assigned to that broader agenda.

The manuscript and exact-arithmetic verification code were generated with
GPT 6 Astra (configured model \texttt{gpt-6-astra}, reasoning effort
\texttt{max}). The exact serving revision was not exposed.
A separate manual review pass checked all assumptions, both directions,
integer-capacity endpoints, empty markets, and the all-allocation obstruction.
The accompanying program checked 342,838 labelled allocations, 7,381 complete
small cost profiles, 1,050 seeded instances, and 10,780 rational interval
boundaries; every check passed. These finite checks supplement the general
proof. This contribution makes no Lean-verification claim.

The public manuscript and reproducibility files are supplied at
\url{https://github.com/jbaelaw/budget-ef1-capacity-frontier/tree/85c02d5aa67cfcbf51163a6a91d5fa08e40e89a0}.
The earlier catalogue proof is preserved and cited. Its unbounded example
and equal-duration result are not claimed as new here.

\begin{thebibliography}{9}
\raggedright
\bibitem{elkind}
E.~Elkind. Fair Division of Chores with Budget Constraints.
Section~4.1 in E.~Markakis, R.~Mehta, and Y.~Zick (eds.),
\emph{Fair Division: Algorithms, Solution Concepts, and Applications},
Dagstuhl Reports 14(9), pp.~145--166, 2025.
\url{https://doi.org/10.4230/DagRep.14.9.145}.
\bibitem{eit}
E.~Elkind, A.~Igarashi, and N.~Teh.
\emph{Fair Division of Chores with Budget Constraints}.
SAGT 2024, LNCS 15156, pp.~55--71, 2024.
\url{https://doi.org/10.1007/978-3-031-71033-9_4}.
Full text: \url{https://arxiv.org/html/2410.23979v1}.
\bibitem{prior}
GPT~6 Astra Ultra. \emph{An unbounded-factor obstruction and an
equal-duration EF1 theorem}. Prior research attempt for D63-448,
Problem Hunting with LLMs, repository snapshot
\texttt{00ced55ba1223fb48603e1f15028d753997ee515}.
\url{https://github.com/mehmetmars7/Erdosproblems-llm-hunter/blob/00ced55ba1223fb48603e1f15028d753997ee515/attacks/open_problems/economics/gpt_6_astra_ultra/D63-448.tex}.
\end{thebibliography}
\end{document}
